\section{Synthesis}
\label{sec:synthesis}

This section draws the per-family results of \cref{sec:crypto}--\cref{sec:retrieval} together
into three views: the three-way trade-off that locates each family
(\cref{ssec:design-space}), per-use-case deployment recommendations (\cref{ssec:deployment}),
and the cross-cutting findings and open problems that no single family section could state on
its own (\cref{ssec:findings}). The headline is that no family dominates:
which scheme is deployable depends on the application, on the constraint it can least afford to
concede.

\subsection{Design Space}
\label{ssec:design-space}

Every family maximizes at most two of the three deployment-readiness criteria and pays on the
third, and the security basis axis of \cref{sec:threat-axes} explains why the third differs so
much. Cryptographic inference
(\cref{sec:crypto}) attains exact fidelity on a provable basis and pays one to three orders of
magnitude in performance. A Tier-3 TEE (\cref{sec:tee}) matches that exact fidelity at
near-plaintext speed, but its basis is the silicon: the guarantee reduces to hardware whose side
channels every threat model excludes, and to the availability of the accelerator that holds it.
Static obfuscation (\cref{sec:obfuscation}) is the cheapest, near-plaintext throughout, but
rests on a heuristic basis and gives up a tunable slice of fidelity. Differential privacy keeps
that near-plaintext cost and adds a formal but probabilistic guarantee, paying again in fidelity
through the budget $\varepsilon$ (\cref{sec:dp}). The hybrid split (\cref{sec:hybrid}) recovers
exact fidelity and a basis stronger than static obfuscation on commodity hardware, and pays
instead in the TEE$\leftrightarrow$GPU serialization that its offloading induces. The criterion a
family concedes is forced by the basis it rests on, not incidental, so no single point on this
surface dominates the rest. Which concession a deployment can absorb is what selects the scheme: a
real-time service is bound by its latency budget, a regulated workload by the security basis it
must be able to prove, a cost-sensitive one by the hardware it already has, and a
quality-critical one by the fidelity it can spare. The next section reads the surface through
these constraints, use case by use case.

\subsection{Deployment Readiness}
\label{ssec:deployment}

\begin{table*}[t!]
  \centering \scriptsize\setlength{\tabcolsep}{4pt}
  \caption{Deployment readiness by family, on the three criteria of \cref{sec:intro}, plus
    \emph{client reliance} (thin = enc/dec a query; stateful = holds keys + light work;
    thick = stores access state and runs a multi-round retrieval/compute loop). Ratings are
    qualitative and indicative, not co-measured; see \cref{tab:performance} for the
    performance basis. \tq{}-marked client cells await the per-scheme review
    (\cref{ssec:deployment}).}
  \label{tab:deployment}
  \begin{tabular*}{\textwidth}{@{\extracolsep{\fill}}llllll@{}}
    \toprule
    Mechanism family & Performance & Fidelity & Threat-model fit & Client & Deploy.\ readiness \\
    \midrule
    Cryptographic (MPC/FHE) & L infer./M retr. & H (exact) & H (often exceeds need) & thin--thick & L--M \\
    TEE (Tier 3)            & H (near-plain.)  & H (exact) & M (HW-trust; side-ch.) & \tq & M (HW-gated) \\
    Static obfuscation      & H (near-plain.)  & H (small loss) & L (fixed; several broken) & \tq & M \\
    Differential privacy    & H                & M (recall at small $\varepsilon$) & M (formal, prob.) & \tq & M \\
    Hybrid split            & M--H (few$\times$) & H (exact) & M--H (Tier-2; per-batch) & \tq & H \\
    \bottomrule
  \end{tabular*}
\end{table*}

The same trade reads differently in each use case, because the constraint a deployment can
least afford to concede changes with the task. We give the recommendation use case by use case.

\emph{Generation} is dominated by the cost of autoregressive decoding, so its binding
constraint is latency. A confidential accelerator (\cref{sec:tee}) is the strongest option for
an operator that has one, running the full model at near-plaintext speed, but the hardware is
scarce and few deployments own it. Where it is absent, static obfuscation runs the unmodified
model on commodity GPUs at near-plaintext cost, at the price of a heuristic basis. The hybrid
split is a weaker fit here than its commodity-hardware story suggests, because its per-token
TEE$\leftrightarrow$GPU round-trips (\cref{sec:hybrid}) make decoding latency its weak point.
MPC and FHE pay one to three orders of magnitude, which is hard to justify for single-user
outsourcing; their natural setting is instead \emph{collaborative} inference among mutually
distrusting data holders, such as joint analysis over disjoint institutional records, where no
party can be trusted to aggregate and a cryptographic guarantee across parties is the whole
point.

\emph{Embedding and reranking} run a smaller, often encoder-sized model, which shifts the
calculus: the workload does not need a flagship accelerator, so paying for a confidential GPU
is overkill. This is the hybrid split's natural setting. It keeps a small trusted core in a
commodity CPU-TEE and offloads to untrusted commodity GPUs (\cref{sec:hybrid}), giving exact
output without enterprise hardware, and the short, non-autoregressive workload makes its
round-trip cost far less acute than in generation. Static obfuscation and local DP are the
near-plaintext alternatives where a heuristic or probabilistic basis is acceptable, and a
cryptographic encoder (SHAFT, Euston) is the choice only when an exact and provable guarantee
is required.

\emph{Vector retrieval} divides by what must be hidden. When the access pattern and the
inter-document structure must stay hidden, the access-hiding schemes answer it: private
information retrieval (Panther, PIR-RAG) returns a record without revealing which one, and
homomorphic scoring (RemoteRAG, TRSE) keeps the similarity values encrypted. When near-plaintext
speed is the priority and a bounded, quantified comparability leak is tolerable, CAPRISE is the
strong choice. Its conditional distance-comparison-preserving encryption hides the embeddings at
roughly $1\times$ cost, its query-side differential privacy blocks embedding inversion, and
unlike SAP it conceals the inter-document structure (\cref{ssec:retr-dense}). The deciding
question is whether that residual leak is acceptable for the corpus at hand.

\emph{Graph retrieval} has no near-plaintext option, so the choice trades cost against which
structural channel it closes (\cref{ssec:retr-graph}). Hiding which node a query touches takes
either client-side cryptography (batched PIR in Pacmann, an ORAM-wrapped index in Compass) or
the hardware route of a TEE (Opal); the cryptographic route shifts cost and state onto the
client, the TEE route onto trusted hardware. Hiding the adjacency volume takes a volume-hiding
encrypted multi-map (XorMM, FLASH) at the cost of storage overhead. Where neither cost is affordable,
anonymization (PrivGemo, ARoG) is the cheapest protection, at a heuristic and unquantified
residual leak. No single scheme closes all three channels at once.

\begin{openproblem}
\label{op:no-e2e}
\emph{No end-to-end confidential pipeline.} Every scheme here protects one or two stages, and no
single design spans embedding, retrieval, and generation; the gaps between stages are uncovered.
Composing per-stage mitigations is largely unstudied: the
joint guarantee of a stacked pipeline, and whether its fidelity losses and residual leakages
compound, are uncharacterized (\cref{op:dp-compose}). The prior RAG-privacy SoK reaches the same gap~\citep{bodea_sok}. \emph{Which compositions hold their per-stage guarantees, and at what
combined fidelity cost?}
\end{openproblem}

\subsection{Cross-Cutting Findings}
\label{ssec:findings}

\paragraph{Client reliance, offline setup, and multi-tenancy.}
The trust-anchor axis of \cref{sec:threat} records where the \emph{trusted} computation runs,
and by construction it brackets the client, since a client need not be named as its own trusted
base. That convention hides how much work several schemes push back onto the client, which this
finding makes explicit. Cryptographic schemes make the user a compute party in 2PC, or drive
private retrieval from the client in PIR and ORAM (\cref{sec:crypto}, \cref{sec:retrieval}).
Covariant weight obfuscation (AloePri, KV-Cloak, ObfusLM) has the client transform the weights
and hold the secret, so the served checkpoint is re-keyed per client and one shared model can no
longer serve many tenants at once. Local DP is client-reliant by construction: the noise must be
added before the data leaves the device, so the client runs the model's input-side layers and
draws the noise itself. DP-Forward makes the cost concrete at the split point, since moving more
layers onto the client raises utility but also its compute and storage~\citep{du_dpforward}; the
offline fine-tuning that calibrates the noise runs on the operator, so the standing burden is
the online forward pass, not the one-time setup. The hybrid split mostly avoids a heavy client
setup because it hides dynamically. SCX is the exception that keeps the client active during
decoding, encoding each new token and recovering the first and last transformer blocks on the
device while the cloud enclave handles the rest, a light but standing per-token role rather than
a one-time cost~\citep{zeng_scx}. Dense and graph retrieval add their own client load: the
client embeds and perturbs its query, Compass keeps an ORAM position map and runs a multi-round
controller, and Pacmann preprocesses the database client-side (\cref{sec:retrieval}). Two
reasons drive all of this. One is structural: local DP and covariant weight obfuscation can
only originate the secret on the client and so forgo multi-tenancy. The other is performance: a
client-side cache, compute party, or embedder lowers server cost or saves rounds. The TEE is the workaround rather than a counterexample: Opal hides the same access
pattern with only a key and a counter on the client (\cref{sec:retrieval}), trading the client's
burden for trust in the hardware.

\paragraph{Confidential accelerators are scarce and enterprise-only.}
The pure confidential-accelerator path, a GPU- or NPU-TEE trusted end to end, is the most
performant route that is also exact and near-plaintext, but it trades the confidentiality
problem for a hardware-procurement one (\cref{sec:tee}). The confidential accelerators our
corpus draws on are the NVIDIA confidential-computing GPUs (the H100/H200/B200 class) and the
Huawei Ascend NPU; the platform space beyond these is narrow and still forming, since confidential GPU
support is recent, concentrated in high-end enterprise parts, and gained inline TEE-I/O only with
Blackwell (\cref{sec:tee}). Trusting such an accelerator end to end also binds a deployment to
one vendor's silicon and attestation. That cost is easy to justify for frontier generation,
where the model is large and the accelerator would be used regardless, and hard to justify for
embedding or reranking, where a small encoder does not need a flagship GPU yet no cheaper
confidential accelerator exists to match the smaller workload. The scarcity is what the hybrid
split sidesteps by offloading to commodity untrusted GPUs (\cref{sec:hybrid}).

\paragraph{For heuristic schemes, a refreshed secret is what survives accumulation.}
The same failure recurs across three families. A single secret, reused on every request, lets
the adversary accumulate the externalities the scheme offloads (the obfuscated activations,
attention, and externalized KV-cache) until they pin down what the secret was meant to hide: the
hidden states, and through them the input. Fixed permutations are broken by prompt reconstruction
in cryptographic inference (\cref{sec:crypto}), static obfuscation accumulates into blind source
separation and learned inversion (\cref{sec:obfuscation}), and a reused KV-cache shuffle is
broken by collision (\cref{sec:hybrid}). In every case the schemes that survive are the ones that
refresh the secret: KV-Cloak draws a fresh permutation per cache block, GELO a fresh mixing
matrix per batch, and SCX a one-time key per generated token. Refreshing is necessary but not
sufficient. Under an orthogonal mixing the Gram matrix $U^{\top}U$ equals $H^{\top}H$ whichever
matrix is drawn, so a fresh secret per batch still exposes the hidden states' pairwise inner
products, which feed the blind-source-separation and anchor attacks of \cref{sec:hybrid}. How
often to refresh under a latency budget is itself open (\cref{op:refresh-frequency}).

\paragraph{Malicious-operator security is an added integrity layer, not a free upgrade.}
Most schemes here assume an honest-but-curious operator that follows the protocol. Defending one
that may also deviate, by dropping, replaying, or tampering, means adding integrity on top of
confidentiality, and that comes by one of two routes. The cryptographic route adds authenticated
computation: information-theoretic MACs make secret-sharing MPC actively secure in the SPDZ line,
TwinShield's U-Verify checks a possibly malicious GPU with a Freivalds-style test
(\cref{sec:hybrid}), and Opal authenticates memory against a malicious host with a Merkle tree
(\cref{sec:tee}). The hardware route adds attestation, as in the confidential accelerators
(Nvidia-CC, Ascend-CC). Malicious resistance is therefore not hardware-exclusive. What stands out
is that the surveyed cryptographic-inference schemes take neither route: SHAFT, Fission, PermLLM,
and CryptoMoE are all proved only semi-honest, and a malicious client can already extract the
model from a semi-honest protocol (\cref{sec:crypto}). Defending against a deviating operator is
narrower than full malicious security, which also covers a cheating client and costs more again;
the surveyed schemes provide neither. MACs could in principle make their linear-algebra core
actively secure at a cost, but the permutation-reveal speed tricks that make these schemes fast
sit outside what a MAC protects, so malicious security there is not a drop-in. A gap remains: malicious resistance is achievable
cryptographically, yet no surveyed crypto scheme pays for it, and the families that do reach it
lean on integrity machinery or trusted hardware that the cheap, semi-honest designs omit.

\paragraph{The field-wide side-channel blind spot.}
\label{ssec:sidechannel}
The residual-leakage axis of the taxonomy (\cref{tab:threat-taxonomy}) gives a uniform answer:
\emph{every} scheme places the microarchitectural, cache, timing, and physical side channels
(power analysis and the memory bus among them) out of scope. GELO excludes accelerator-side leakage beyond raw memory
reads~\citep{belikov_gelo}, Opal enumerates the in-enclave channels it does \emph{not}
defend~\citep{kaviani_opal}, and the rest omit them. This is a deliberate exclusion, and it
falls hardest on the TEE and hybrid families, whose entire trust is the silicon. Their
best-documented weaknesses are exactly microarchitectural and physical: transient-execution
attacks of the Spectre and Foreshadow class~\citep{spectre,foreshadow}, and physical memory-bus
interposition such as TEE.Fail, which extracts keys and forges attestation from the
deterministic memory encryption of server TEEs. The vendors classify such physical attacks as
out of scope, with no fix planned~\citep{tee_fail}. Individual attacks are patched on the platform cycle, but
the class recurs because each generation's performance optimizations reopen an observable
channel, and unlike a cryptographic assumption a hardware-anchored design cannot close it
without the vendor.

\paragraph{Graph-RAG is the most privacy-underserved setting.}
No surveyed work delivers an end-to-end confidential graph-RAG deployment, and it is the setting
the field, ours included, covers least. What this SoK offers is a per-component map of which
surveyed primitive serves which part of the graph pipeline (\cref{ssec:retr-graph}): entity and
relation ANN lookups by Compass, Pacmann, and CAPRISE; one-hop adjacency by the volume-hiding
multi-maps XorMM and FLASH; graph storage by PeGraph and GORAM; access hiding by Opal (TEE) or a client-side ORAM index;
and the cheapest, heuristic protection by the anonymization of PrivGemo and ARoG. One stage
stays uncovered across the corpus: building the graph index itself under confidentiality, before
any query runs. Confidential graph construction is the clearest open problem in this setting.

% tab:gapmap rows = pipeline stages; columns = mechanism families (§3-§7);
% cells = representative schemes; ``---'' = no practical scheme (a gap).
%\begin{table*}[t!]
%  \centering \scriptsize
%  \caption{Coverage of the inference-and-RAG pipeline by mechanism family. Cell = one
%    or two representative schemes (\cref{sec:crypto}--\cref{sec:hybrid}; see those
%    tables for the full set); ``---'' = no practical scheme known to us (a gap).}
%  \label{tab:gapmap}
%  \begin{tabular*}{\textwidth}{@{\extracolsep{\fill}}llllll@{}}
%    \toprule
%    Pipeline stage & Crypto/MPC & TEE & Obfuscation & DP & Hybrid split \\
%    \midrule
%    Embed            & SHAFT/Euston       & H100-CC          & SGT/Eguard    & DP-Fwd/SPARSE & GELO/ObfuscaTune \\
%    Store (vectors)  & CAPRISE/SAP        & Opal             & ---           & ---           & --- \\
%    Retrieve (ANN)   & Compass/Pacmann    & Opal+ORAM        & ---           & RemoteRAG     & --- \\
%    Rerank           & p\textsuperscript{2}RAG/TRSE & Opal   & ---           & RemoteRAG(HE) & --- \\
%    Generate         & PermLLM/CryptoMoE  & PipeLLM          & AloePri/OSNIP & DP-RAG(out)   & GELO/SCX \\
%    Graph construct. & ---                & ---              & ARoG          & ---           & PrivGemo \\
%    Graph traverse   & Pacmann/XorMM      & Opal/H$_2$O$_2$RAM & ARoG        & ---           & PrivGemo \\
%    \bottomrule
%  \end{tabular*}
%\end{table*}

%\begin{table*}[t!]
%  \centering \footnotesize
%  \caption{Performance and fidelity by use case (indicative; normalized per use case,
%    \cref{sec:comparison}). ``$\times$'' = overhead vs.\ plaintext; ``exact'' = computes the
%    true function. Cross-paper figures are indicative, not co-measured. All figures are
%    \emph{paper-reported}; we do not independently reproduce them, and Chen et
%    al.~\citep{chen_sok_pti} find that re-benchmarked crypto-inference numbers can diverge from
%    the originals---so cross-scheme readings here are order-of-magnitude, and open-source
%    availability is tracked per scheme in the family tables (\cref{sec:crypto}--\cref{sec:hybrid}).}
%  \label{tab:performance}
%  \begin{tabular*}{\textwidth}{@{\extracolsep{\fill}}>{\raggedright\arraybackslash}p{1.6cm}>{\raggedright\arraybackslash}p{6.3cm}>{\raggedright\arraybackslash}p{3.6cm}>{\raggedright\arraybackslash}p{3.6cm}@{}}
%    \toprule
%    Use case & Mechanisms (representative) & Overhead (normalized) & Fidelity \\
%    \midrule
%    Generation & MPC (PermLLM, CryptoMoE, Fission); TEE (PipeLLM); obfusc (AloePri, SGT, OSNIP); hybrid (GELO, SCX, ObfuscaTune)
%               & MPC 1--2 orders; TEE $<$20\%; obfusc $\sim$0; hybrid few$\times$
%               & exact (MPC/TEE/hybrid); lossy (obfusc) \\
%    Embedding (encoder) & MPC (SHAFT); FHE (Euston); DP (DP-Forward, SPARSE); obfusc (Eguard)
%               & crypto 1--2 orders; obfusc/DP $\sim$0
%               & exact (crypto); acc.\ $\Delta$ (obfusc/DP) \\
%    Vector retrieval & crypto (p\textsuperscript{2}RAG, Panther, PIR-RAG, RAGtime-PIANO); DP (RemoteRAG); enc (CAPRISE, SAP)
%               & sub-s--10\,s @$10^6$--$10^7$; enc $\sim$0
%               & exact (PIR/FHE); recall $\Delta$ (DP); order-leak (DCPE) \\
%    Graph retrieval & TEE (Opal); crypto (Compass, Pacmann); SSE (XorMM, PeGraph, GORAM); anon (ARoG, PrivGemo)
%               & graph-ANN sub-s--s; KG-anon $\sim$0
%               & exact (ORAM/SSE); quality-only (anon) \\
%    \bottomrule
%  \end{tabular*}
%\end{table*}

